% Wave-17 B2/20: NCCR quiver-chart covers of compact CY_3 and the
% assembly of chart-level positive halves Y^+_U into a global Y^+(X).
% Voice: Van den Bergh (tilting/MCM/Auslander) + Bocklandt (moduli-
% theoretic quiver-with-potential). Chriss--Ginzburg discipline.
% Primary sources (all cross-checked):
%   Van den Bergh, M., "Non-commutative crepant resolutions",
%     in The Legacy of Niels Henrik Abel (Oslo 2002), Springer 2004,
%     pp.~749--770; also arXiv:math/0211064.
%   Van den Bergh, M., "Three-dimensional flops and non-commutative
%     rings", Duke Math.~J.~122 (2004) 423--455 (arXiv:math/0207170).
%   Bocklandt, R., "Generating toric noncommutative crepant
%     resolutions", J.~Algebra 364 (2012) 119--147 (arXiv:1104.1597);
%     Bocklandt, R., "Calabi--Yau algebras and quiver polyhedra",
%     Math.~Z.~273 (2013) 311--329 (arXiv:0905.0232).
%   Kontsevich, M., Soibelman, Y., "Stability structures, motivic
%     Donaldson--Thomas invariants and cluster transformations",
%     arXiv:0811.2435 (especially Section~7, Hall-algebra categorification).
%   Szendroi, B., "Non-commutative Donaldson--Thomas theory and the
%     conifold", Geom.~Topol.~12 (2008) 1171--1202 (arXiv:0705.3419).

\providecommand{\Hom}{\mathrm{Hom}}
\providecommand{\End}{\mathrm{End}}
\providecommand{\Ext}{\mathrm{Ext}}
\providecommand{\Spec}{\mathrm{Spec}}
\providecommand{\Coh}{\mathrm{Coh}}
\providecommand{\Mod}{\text{-}\mathrm{Mod}}
\providecommand{\modd}{\text{-}\mathrm{mod}}
\providecommand{\Db}{D^{b}}
\providecommand{\Perf}{\mathrm{Perf}}
\providecommand{\Rep}{\mathrm{Rep}}
\providecommand{\CoHA}{\mathrm{CoHA}}
\providecommand{\Jac}{\mathrm{Jac}}
\providecommand{\Crit}{\mathrm{Crit}}
\providecommand{\Eone}{E_{1}}
\providecommand{\cA}{\mathcal A}
\providecommand{\cH}{\mathcal H}
\providecommand{\cO}{\mathcal O}
\providecommand{\cT}{\mathcal T}
\providecommand{\cU}{\mathcal U}
\providecommand{\fU}{\mathfrak U}
\providecommand{\bC}{\mathbb C}
\providecommand{\bZ}{\mathbb Z}

\section*{Wave-17 B2: NCCR quiver-chart covers of compact CY$_3$}

\subsection*{Cycle 1. Local NCCR data per affine chart.}

Let $X$ be a compact smooth Calabi--Yau threefold, $\fU=\{U_{\alpha}\}$
a Zariski cover by affines, each $U_{\alpha}=\Spec R_{\alpha}$ Gorenstein
and admitting a crepant (possibly non-commutative) desingularisation.
Van den Bergh 2004 (\emph{Non-commutative crepant resolutions}, Def.~4.1)
attaches to each such $R_{\alpha}$ an NCCR
\[
  \Lambda_{\alpha}\;:=\;\End_{R_{\alpha}}(T_{\alpha}),
  \qquad T_{\alpha}\in\mathrm{refl}(R_{\alpha}),\;
  \mathrm{gldim}\,\Lambda_{\alpha}<\infty,\;
  \Lambda_{\alpha}\in\mathrm{CM}(R_{\alpha}).
\]
When $U_{\alpha}$ is a small open containing a compact curve
$C_{\alpha}\subset X$ with normal bundle of the Bridgeland type, Van den
Bergh 2004 (\emph{Three-dimensional flops}, Thm.~B, arXiv:math/0207170)
produces a tilting bundle $\cT_{\alpha}\in\Db\Coh(U_{\alpha})$ with
$\End(\cT_{\alpha})=\Lambda_{\alpha}$, and a derived equivalence
$\Db\Coh(U_{\alpha})\simeq\Db(\Lambda_{\alpha}\modd)$. By Bocklandt 2013
(\emph{Calabi--Yau algebras and quiver polyhedra}, Thm.~3.3;
arXiv:0905.0232), $\Lambda_{\alpha}$ is presented as a Jacobi algebra
$\Jac(Q_{\alpha},W_{\alpha})=\bC Q_{\alpha}/\langle\partial_{a}W_{\alpha}
\rangle_{a\in Q_{\alpha,1}}$ of a quiver with potential, and this
presentation is unique up to right-equivalence of potentials. The key
CY$_{3}$-algebra property --- bimodule $3$-Calabi--Yau in the sense of
Ginzburg --- is built in: the Koszul-type resolution of the diagonal
$\Lambda_{\alpha}$-bimodule
\[
  0\to\Lambda_{\alpha}\otimes\Lambda_{\alpha}
   \to\Lambda_{\alpha}\otimes V_{\alpha}\otimes\Lambda_{\alpha}
   \to\Lambda_{\alpha}\otimes V_{\alpha}\otimes\Lambda_{\alpha}
   \to\Lambda_{\alpha}\otimes\Lambda_{\alpha}\to\Lambda_{\alpha}\to 0,
\]
with $V_{\alpha}=\mathrm{span}(Q_{\alpha,1})$ (Ginzburg 2006, Prop.~5.1.9;
internal to Bocklandt's treatment) exhibits
$\Lambda_{\alpha}^{!}\cong\Lambda_{\alpha}[3]$, i.e.\ the Verdier dual is
a shifted copy.

\subsection*{Cycle 2. Chart-level positive halves
  $Y^{+}_{\alpha}=\CoHA(\Lambda_{\alpha},W_{\alpha})$.}

The critical cohomological Hall algebra of Kontsevich--Soibelman 2008
(arXiv:0811.2435, \S 7) attaches to $(Q_{\alpha},W_{\alpha})$ the graded
associative algebra
\[
  \cH(Q_{\alpha},W_{\alpha})
  \;=\;\bigoplus_{\mathbf d\in\bZ^{Q_{\alpha,0}}_{\geq 0}}
    H^{\mathrm{BM}}_{*}\!\bigl(\Crit(W_{\alpha,\mathbf d});
      \phi_{W_{\alpha,\mathbf d}}\bigr),
\]
with product the proper pushforward along the stack of short exact
sequences of $\Lambda_{\alpha}$-modules. Schiffmann--Vasserot 2013
($\bC^{3}$; arXiv:1202.2756) and Rapcak--Soibelman--Yang--Zhao 2020
(toric CY$_{3}$ without compact $4$-cycles; arXiv:2007.13365) identify
$\cH(Q_{\alpha},W_{\alpha})$ with the positive half
$Y^{+}(\widehat{\mathfrak g}_{Q_{\alpha}})$ of the affine super-Yangian
attached to $Q_{\alpha}$: thus $Y^{+}_{\alpha}:=Y^{+}
(\widehat{\mathfrak g}_{Q_{\alpha}})$ is the chart-level invariant. Its
Poincar\'e series equals the Kontsevich--Soibelman DT generating series
of $(Q_{\alpha},W_{\alpha})$ (KS 2008, Thm.~8). This is the chart-level
datum from which a putative global $Y^{+}(X)$ must be assembled.

\subsection*{Cycle 3. Overlap compatibility: Morita, not isomorphism.}

On $U_{\alpha\beta}=U_{\alpha}\cap U_{\beta}$, \emph{two} NCCR
presentations co-exist: $\Lambda_{\alpha}|_{U_{\alpha\beta}}$ and
$\Lambda_{\beta}|_{U_{\alpha\beta}}$. Van den Bergh's uniqueness-up-to-MCM
theorem (2004 \emph{Three-dimensional flops}, Cor.~4.5) does \emph{not}
assert $\Lambda_{\alpha}\cong\Lambda_{\beta}$ on overlap; only that the
two NCCRs are derived-equivalent via a tilting bimodule
\[
  {}_{\Lambda_{\alpha}}\cT_{\alpha\beta}{}_{\Lambda_{\beta}}
  \;\in\;\Db(\Lambda_{\alpha}\otimes\Lambda_{\beta}^{\mathrm{op}}),
  \qquad
  \cT_{\alpha\beta}\otimes^{L}_{\Lambda_{\beta}}(-)\colon
  \Db(\Lambda_{\beta}\modd)\xrightarrow{\sim}\Db(\Lambda_{\alpha}\modd).
\]
The conifold case makes this concrete: the two small crepant resolutions
$\cY_{\pm}$ furnish tilting bundles $\cO\oplus\cO(\pm 1)$ whose
endomorphism algebras are the \emph{same} $\Lambda_{\mathrm{NCCR}}=
\End_{R}(R\oplus I)$ (Szendroi 2008, Prop.~2.5); but for a general pair
of flopping small resolutions the endomorphism algebras differ, and only
Morita-equivalence persists. Bocklandt 2012 (\emph{Generating toric
NCCRs}, Thm.~1.2; arXiv:1104.1597) gives an explicit recipe for the
tilting bimodule $\cT_{\alpha\beta}$ in the toric setting as a
right-equivalence of quivers-with-potential, together with a combinatorial
criterion (``consistency'' of the dimer / brane-tiling) that lets two
toric charts glue without obstruction.

\subsection*{Cycle 4. Global derived stack of quiver representations.}

The correct global object is \emph{not} a single quiver with potential on
$X$; no such object exists in general. It is the stack
\[
  \mathfrak M_{X}\;:=\;\operatorname*{hocolim}_{\alpha\in N(\fU)}
    \Rep(\Lambda_{\alpha})\;/\;\text{gauge},
\]
where $N(\fU)$ is the nerve of the cover, each $\Rep(\Lambda_{\alpha})=
\bigsqcup_{\mathbf d}\Rep_{\mathbf d}(\Lambda_{\alpha})$ is the Bocklandt
moduli stack of $\Lambda_{\alpha}$-representations (Bocklandt 2013,
\S 4), and the transition maps on $U_{\alpha\beta}$ are the Morita
equivalences of Cycle 3 realised as stack equivalences
\[
  \cT_{\alpha\beta}\otimes_{\Lambda_{\beta}}^{L}(-)
  \colon\Rep(\Lambda_{\beta})\bigr|_{U_{\alpha\beta}}
  \xrightarrow{\sim}\Rep(\Lambda_{\alpha})\bigr|_{U_{\alpha\beta}}.
\]
Kontsevich--Soibelman 2008 (\S 7.8) observes that such a hocolim, when
the cocycle $\{\cT_{\alpha\beta}\}$ satisfies the evident $2$-cocycle
condition $\cT_{\alpha\beta}\otimes\cT_{\beta\gamma}\simeq
\cT_{\alpha\gamma}$ on triple overlaps, is a derived Artin $1$-stack;
its tangent complex at a point $E\in\Coh(X)$ is $\Ext^{*}_{X}(E,E)[1]$,
matching the chart-level $\Ext^{*}_{\Lambda_{\alpha}}$ under the Morita
equivalence of Cycle 3. (Bocklandt 2013 Thm.~1.1 verifies the
$2$-cocycle condition for toric covers; in the non-toric compact CY$_{3}$
setting the condition is a separate hypothesis, not a theorem.)

\subsection*{Cycle 5. Global positive half as sheafy derived sections.}

Define the presheaf of CoHAs on $\fU$ by
$\underline{Y^{+}}(U_{\alpha})=Y^{+}_{\alpha}$ with restriction maps
$\underline{Y^{+}}(U_{\alpha})\to\underline{Y^{+}}(U_{\alpha\beta})$ the
pullback along the Morita equivalence $\cT_{\alpha\beta}$. Pass to the
associated sheaf of $\Eone$-algebras $\underline{Y^{+}}^{\sharp}$ on $X$;
the global positive half is
\[
  Y^{+}(X)\;:=\;R\Gamma\bigl(X,\underline{Y^{+}}^{\sharp}\bigr).
\]
Two finite-dimensional cross-checks:
(i) For $X=\bC^{3}/\Gamma$ with $\Gamma\subset\mathrm{SU}(3)$ finite
(an NCCR via the skew group algebra $\bC[x,y,z]\#\Gamma$ after Bridgeland--
King--Reid), the cover is a single chart and $Y^{+}(X)=Y^{+}(\widehat{\mathfrak
g}_{Q_{\Gamma}})$ recovers the McKay-quiver positive half.
(ii) For $X$ a resolved toric CY$_{3}$ without compact $4$-cycles,
Rapcak--Soibelman--Yang--Zhao 2020 prove
$Y^{+}(X)\cong Y^{+}(\widehat{\mathfrak g}_{Q_{X}})$ and the nerve of
$\fU$ has trivial higher cohomology (toric-Cech acyclicity), so the
hocolim collapses to the single toric quiver. For compact non-toric $X$
(quintic; $K3\times E$) the hocolim is \emph{not} known to collapse and
the status of $Y^{+}(X)$ is: conjectural, with the derived global-sections
formula above as the proposed definition. The $\kappa_{\mathrm{ch}}$
invariants are chart-additive along the \v{C}ech spectral sequence
$E_{2}^{p,q}=\check H^{p}(\fU,\underline{\kappa_{\mathrm{ch}}})\Rightarrow
\kappa_{\mathrm{ch}}(X)$, with the chart-level
$\kappa_{\mathrm{ch}}(\Lambda_{\alpha})$ given by the Bocklandt 2013
\S 5 character formula for Jacobi algebras; the $K3\times E$ value
$\kappa_{\mathrm{ch}}=0$ of the cache matches $\chi(\cO_{K3\times E})=0$
by Hodge-supertrace K\"unneth.

\subsection*{Status of the five cycles.}

Cycle 1 (local NCCR) and Cycle 2 (chart-level $Y^{+}_{\alpha}$) are
unconditional per the cited Van den Bergh, Bocklandt, and Kontsevich--
Soibelman / Schiffmann--Vasserot / Rapcak--Soibelman--Yang--Zhao results.
Cycle 3 (Morita overlap) is established in the toric setting (Bocklandt
2012 Thm.~1.2) and conjectural in the non-toric compact case. Cycle 4
(global derived stack) is unconditional once the $2$-cocycle condition
holds; verifying it on a specific compact non-toric $X$ is an open problem.
Cycle 5 (sheafy $Y^{+}(X)$) is therefore \emph{conjectural} as a
definition for compact non-toric CY$_{3}$, and cohomologically non-trivial
only when the cover's nerve has higher cohomology --- the obstruction
to a single global NCCR.
